\section{Main theorem suite and proof architecture}
\label{sec:block14-main-suite}

Let
\[
 \mathfrak A_{\mathrm{cert}}:=\mathfrak A_{13}
\]
denote the finite certified atlas constructed in \cref{sec:block13-global-invariant}.  Its reversible core is the disjoint union of five connected semantic components, while its directional satellites retain explicit loss or obstruction passports.

\begin{theorem}[Stratified resolved representation, descent, and complete invariant theorem]
\label{thm:block14-main}
Let \(\mathbb K\) be the resolved theory used in this manuscript and let \(\mathfrak A_{\mathrm{cert}}\) be the declared finite atlas.  Then the following statements hold.
\begin{enumerate}
\item Every chart is a profile-preserving realization of the common resolved syntax, and every reversible edge is an equivalence on its named overlap category.
\item The descent category exists as the finite pseudolimit of the atlas and decomposes as the coproduct of its five connected semantic components.
\item Every proof whose constructor support lies in the common verified preservation profile transports to compatible local theorems and hence to the corresponding global resolved object.
\item The semantic normalizers induce the comparison 2-cells on all nonempty triple overlaps and force the quadruple-overlap coherence equation.
\item The compatible native invariant family descends to a tagged global invariant
\[
 \mathcal J_{\mathrm{cert}}:
 \Desc(\mathfrak A_{\mathrm{cert}})
 \longrightarrow
 \Desc(\mathcal V_{\mathrm{cert}}).
\]
\item The descended invariant is complete on the declared normalized global category:
\[
 \mathcal J_{\mathrm{cert}}(X)\cong
 \mathcal J_{\mathrm{cert}}(Y)
 \quad\Longleftrightarrow\quad
 X\cong Y.
\]
\item Every noninvertible construction remains outside the reversible core and carries a typed record of its preserved structure, forgotten information, named residual, or obstruction class.
\end{enumerate}
\end{theorem}

\begin{proof}
The common resolved syntax and preservation-profile transport are established in \cref{sec:resolved-foundations}.  The finite atlas and its pseudolimit are supplied by \cref{thm:atlas-assembly}.  The normalizer comparison cells and their higher coherence are \cref{thm:block13-normalizer-coherence}.  Descent of compatible invariant families is \cref{thm:block13-certified-invariant-descent}.  Completeness is \cref{thm:block13-global-completeness}; its hypotheses are verified component by component in \cref{sec:block13-global-invariant}.  The bridge passports and obstruction taxonomy in \cref{sec:strata-bridges-obstructions} keep every directional construction outside the maximal subgroupoid.  Combining these results proves all seven assertions.
\end{proof}

\begin{corollary}[Relative universality]
\label{cor:block14-relative-universality}
The certified atlas is universal for profile-preserving models of the declared resolved grammar on its five normalized semantic components.  This universality is relative to the chosen grammar, strata, bridges, preservation profiles, normalization targets, and descent doctrine.
\end{corollary}

\begin{boundary}[Publication boundary]
\label{bnd:block14-publication-boundary}
\Cref{thm:block14-main} does not imply a global canonical normal form for arbitrary mathematical objects, equivalence of all broad strata, decidability of arbitrary presentation equivalence, or effectivity of an unrestricted heterogeneous stack.  It also does not identify functional-Mahler equations with Newton-Mahler interpolation, forget branch data in the algebraic-to-holonomic direction, or infer Tannakian symmetry from a merely faithful additive realization.
\end{boundary}

\begin{table}[htbp]
\centering
\small
\caption{The publication theorem spine.  Section numbers are resolved by the final build; the logical order, rather than the numbering, is authoritative.}
\label{tab:block14-theorem-suite}
\begin{tabularx}{\textwidth}{@{}L{0.16\textwidth}L{0.22\textwidth}YL{0.14\textwidth}@{}}
\toprule
Stage & Principal result & Exact role in the proof & Location \\
\midrule
Resolved source & Additive-Karoubi envelope & Literal finite packet objects and profile-indexed realization source & \cref{thm:additive-karoubi-envelope} \\
Typed atlas & Bridge passports and obstruction separation & Reversible overlaps are separated from directional loss & \cref{sec:strata-bridges-obstructions} \\
Finite assembly & Profiled atlas assembly & Bilimit, theorem transport, groupoid core, and global directional functors & \cref{thm:atlas-assembly} \\
Native normalization & Certified stratum atlases & Complete local coordinates on declared semantic slices & Part II \\
Higher coherence & Normalizer-induced coherence & Unique semantic normalizers force all certified overlap cycles to commute & \cref{thm:block13-normalizer-coherence} \\
Invariant descent & Certified invariant-descent theorem & Coherent local invariants descend through effective object and morphism gluing & \cref{thm:block13-certified-invariant-descent} \\
Global completeness & Tagged componentwise invariant & The descended invariant reflects isomorphism on the five-component union & \cref{thm:block13-global-completeness,thm:block14-main} \\
Lane transport & Finite profiled Universal Lane Transport Theorem & Exact block transport, refinement, selection, residuals, conditional fold loss, and certified descent & \cref{thm:block15-universal-lane-transport,thm:block15-lane-descent} \\
Twisted transport & Holonomy-Twisted Universal Lane Transport Theorem & Matrix-local-system synthesis, base-fiber invariant normal form, orbit coarsening, and orbital path algebra & \cref{thm:block16-holonomy-twisted-lane-transport,cor:block16-base-fiber-invariant-normal-form,thm:block16-orbit-descent,thm:block16-permutation-orbital-calculus} \\
\bottomrule
\end{tabularx}
\end{table}


\begin{figure}[htbp]
\centering
\includegraphics[width=.97\textwidth]{publication_proof_architecture.pdf}
\caption{Proof architecture of the journal-facing theorem.  The vertical spine contains the reversible construction and descent argument.  Directional satellites and obstruction data remain typed external inputs rather than silent equivalences.}
\label{fig:block14-publication-architecture}
\end{figure}
